\section{Scalar descent through quasi-\'etale uniformization}
\label{sec:r24-harmonic-uniformization}

\paragraph{Status and authorship.}
The scalar assertions proved below are \texttt{DIRECT\_COROLLARY} of the
specified primary results.  They are not an arbitrary-current descent
theorem and do not settle KC12.  The independent-audit agent constructed
this deduction after the root requested a source audit.  The root
independently proposed the same Galois-saturation repair after being
informed of the relevant passage in Ducros.  Neither participant is an
author-independent reviewer of this deduction.  A separate review is
required before it is used as an accepted input.

\subsection{Definitions and exact external interfaces}

Let $k$ be a complete nontrivially valued non-Archimedean field.  In its
application to this project, $k$ is the original $K$, or a finite
extension used only in scalar analytic geometry, with the given
absolute value extended and never renormalized.  All spaces in the
descent assertion are good, Hausdorff and strictly $k$-analytic.

The notation $\operatorname{PL}_{\mathbf R}(X)$ has the geometric
meaning of Gubler--Jell--Rabinoff, Definition~5.1: on a covering by
strictly analytic domains the function is a finite real linear
combination of logarithms of absolute values of analytic units on
those domains.  This definition concerns functions on the full
analytic domains.  It is not a definition in terms of a selected
skeleton or a chosen finite set of ordinary current tests.  We use
the following primary inputs, with their original hypotheses.

\begin{enumerate}
\item Gubler--Jell--Rabinoff, arXiv:2111.05741v3, Definition~5.1,
Remark~5.2 and Lemma~5.4: geometric real PL functions are G-local and
are preserved by analytic pullback.  Proposition~5.11, printed
pp.~34--35, gives descent along finite, flat, surjective morphisms of
good strictly analytic spaces.  Corollary~5.12, printed p.~36, gives
descent after extension to the completion of an algebraic extension.
\item Ducros, \emph{Families of Berkovich spaces}, Theorem~5.4.6 and
the reduction in the proof of Proposition~9.1.1(c), printed
pp.~118--120 and~177, give the following local form.  A quasi-\'etale
morphism can be covered G-locally on its source by strict affinoid
domains $D$ admitting a factorization
\[
 D\lhook\joinrel\longrightarrow E\longrightarrow Z
 \lhook\joinrel\longrightarrow X,                         \tag{HU1}
\]
where $D$ is an analytic domain of $E$, $Z$ is a connected strict
affinoid domain of $X$, and $E\to Z$ is connected finite \'etale.
One may use finitely many such pieces when the source is compact.
The same proof explicitly passes to a connected finite Galois cover
of $Z$ dominating $E$.
\item Ducros, Proposition~9.1.1, printed pp.~177--178: if $Y$ is
quasi-compact and $\Gamma$-strict, $X$ is separated, and a coherent
sheaf on $Y$ is $X$-flat with quasi-finite support, its image is a
compact $\Gamma$-strict analytic domain of $X$.  We use only the case
of the structure sheaf of a strict affinoid domain of a finite
\'etale cover, with $\Gamma=\{1\}$.  Such a domain maps
quasi-\'etale, and hence flat, to the base.
\item Temkin, arXiv:1503.05729v2, Theorem~3.3.1 and
Corollary~3.3.2, printed p.~12, together with the rigid/Berkovich
translation in \S3.4: a compact quasi-smooth strictly $k$-analytic
space has, after a finite extension $l/k$, a surjective
quasi-\'etale affinoid cover possessing a \emph{strictly semistable}
formal $l^\circ$-model.  It is this statement, retaining the chosen
strictly semistable model, that is used below.  The separate wording
of Theorem~3.4.1 says only that the maximal affine model is
semistable; that wording alone is not silently strengthened.
\item Gubler--Rabinoff, arXiv:2506.13548v2,
Definition~7.1 and Proposition~7.2(6),(7), printed pp.~34--35:
classically pluriharmonic scalars are preserved by base extension and
analytic pullback.  Proposition~7.8, printed p.~36: on a good,
quasi-compact strictly analytic space having a strictly semistable
formal model, every classically pluriharmonic scalar is geometric
real PL.
\end{enumerate}

Here \emph{classically pluriharmonic} means that both $h$ and $-h$
are classically psh in the curve-test sense of the last reference.
No identification with an arbitrary continuous ordinary-current
harmonic witness is included in that definition.

\subsection{Galois saturation and scalar PL descent}

\paragraph{Proposition (compact quasi-\'etale scalar descent).}
Let $f:Y\to X$ be a surjective quasi-\'etale morphism between compact
good Hausdorff strictly $k$-analytic spaces.  For every real function
$h$ on $X$,
\[
 h\circ f\in\operatorname{PL}_{\mathbf R}(Y)
 \quad\Longleftrightarrow\quad
 h\in\operatorname{PL}_{\mathbf R}(X).                   \tag{HU2}
\]
The assertion is about the given function on all analytic points.
It does not assert that $f$ has an interior lift over each point of
$X$, nor that $f$ becomes finite on an ordinary neighborhood of
every target point.

\paragraph{Proof.}
Analytic pullback proves the reverse implication.  Assume the left
side.  By HU1 and compactness, choose finitely many strict affinoid
source domains $D_i$ covering $Y$, with factorizations
$D_i\hookrightarrow E_i\to Z_i\hookrightarrow X$ as in HU1.
The finite family can be obtained from finitely many source
neighborhoods covered by the local affinoid pieces; thus it is an
actual G-cover and not just a collection of selected points.
It suffices to prove the assertion on the image of each $D_i$.
Fix one piece and omit its index.

Choose a connected finite Galois cover $q:\widetilde E\to Z$
dominating $E\to Z$, and write $G=\operatorname{Aut}_Z(\widetilde E)$.
Let $\widetilde D$ be the inverse image of $D$.  It is a compact
strict affinoid domain of $\widetilde E$.  The composite
$\widetilde D\to Z$ is quasi-\'etale: analytic-domain embeddings and
finite \'etale maps are quasi-\'etale, and this property is stable
under composition.  In particular it is flat and has
zero-dimensional fibers.  Ducros, Proposition~9.1.1, applied to
$\mathcal O_{\widetilde D}$ therefore shows that
\[
 C=q(\widetilde D)=f(D)
 \quad\hbox{is a compact strict analytic domain of }Z.    \tag{HU3}
\]
It is consequently also an analytic domain of $X$.

For every analytic point $z\in Z$, the Galois group acts transitively
on the points of $q^{-1}(z)$.  This is the usual finite-Galois-cover
fiber property, and is also the property used in the Galois
saturation in Ducros's proof of Proposition~9.1.1(c).  More
explicitly, after tensoring the finite \'etale fiber algebra with
$\mathscr H(z)$, its geometric points form a $G$-torsor.  The
absolute Galois group of $\mathscr H(z)$ commutes with this action;
therefore $G$ is transitive on its orbits, equivalently on the field
factors of that finite \'etale algebra.  Each finite extension of the
complete valued field $\mathscr H(z)$ has its unique extended
absolute value, so these factors are exactly the analytic fiber
points.  Hence
\[
 q^{-1}(C)=\bigcup_{g\in G}g(\widetilde D).              \tag{HU4}
\]
Indeed, any point on the right maps to the image of $\widetilde D$;
conversely, a point above $z\in C$ is a conjugate of some chosen
point of $\widetilde D$ above $z$.  Both sides carry their induced
analytic-domain structure in $\widetilde E$, and analytic-domain
structures on a fixed subset are unique.  The union in HU4 is
finite.  No points lying on domain boundaries have been removed.

Put $\widetilde h=h\circ q$ on $\widetilde E$, where $h$ is restricted
to $Z$.  Its restriction to $\widetilde D$ is geometric real PL by
pullback from $D$.  Since $q\circ g=q$, the function $\widetilde h$
is $G$-invariant.  Transporting the presentation on $\widetilde D$
by each $g$ proves that $\widetilde h$ is geometric real PL on every
$g(\widetilde D)$.

For clarity, no goodness assumption is needed for the possibly
non-affinoid union in HU4.  Cover $C$ by strict affinoid domains $V$.
For each such $V$, the space $q^{-1}(V)$ is strict affinoid and
$q^{-1}(V)\to V$ is finite \'etale and surjective.  Its intersections
with the finitely many domains $g(\widetilde D)$ form a finite
analytic-domain cover.  Restriction of the preceding PL
presentations and G-locality show that
$\widetilde h|_{q^{-1}(V)}$ is real PL.  Gubler--Jell--Rabinoff,
Proposition~5.11, now gives $h|_V\in\operatorname{PL}_{\mathbf R}(V)$.
G-locality gives $h|_C\in\operatorname{PL}_{\mathbf R}(C)$, understood
through this affinoid cover if necessary.

Finally, the finitely many compact analytic domains $f(D_i)$ cover
$X$, because the $D_i$ cover $Y$ and $f$ is surjective.  They form a
finite G-cover of $X$.  Applying G-locality once more proves HU2.
\hfill$\square$

\paragraph{What the saturation accomplishes.}
The finite \'etale map used for descent is
$q^{-1}(V)\to V$, after $V$ has been recognized as an actual target
analytic domain.  It need not map back into the \emph{single}
original source piece $D$.  The equality of the scalar pullback on
all its Galois conjugates supplies PL presentations on the additional
pieces.  Thus this proof is not a domination of the original
quasi-\'etale cover by an ordinary finite \'etale cover, and it makes
no assertion of descent for an arbitrary current originally given
only on $D$.

\subsection{The scalar consequence of altered uniformization}

\paragraph{Corollary (classical pluriharmonic scalars are geometric
real PL in the compact strict quasi-smooth case).}
Let $X$ be a compact good Hausdorff quasi-smooth strictly
$k$-analytic space.  Every classically pluriharmonic real function
on $X$ belongs to $\operatorname{PL}_{\mathbf R}(X)$:
\[
 \operatorname{PH}_{\mathrm{cl}}(X)
       \subseteq\operatorname{PL}_{\mathbf R}(X).       \tag{HU5}
\]
In particular this holds over the original $K$.  It also yields the
corresponding local assertion at a point admitting a compact strict
quasi-smooth analytic neighborhood on which the scalar is
classically pluriharmonic.

\paragraph{Proof.}
Use Temkin, Corollary~3.3.2 and \S3.4, to choose a finite extension
$l/k$ and a surjective quasi-\'etale morphism
$p:Y\to X_l$ with $Y$ affinoid and possessing a strictly semistable
formal model over $l^\circ$.  The original absolute value is merely
extended to $l$.  The scalar
$h_l\circ p$ is classically pluriharmonic by
Gubler--Rabinoff, Proposition~7.2(6),(7).  Their Proposition~7.8
shows that it is geometric real PL on $Y$.  Apply HU2 over $l$ to
obtain $h_l\in\operatorname{PL}_{\mathbf R}(X_l)$.  Finally,
Gubler--Jell--Rabinoff, Corollary~5.12, descends this property along
the finite algebraic extension $l/k$ and gives HU5.
All steps concern scalars and their analytic pullbacks.  No current
over $l$ is constructed or imported. \hfill$\square$

\subsection{Boundary pressure test and the remaining KC interfaces}

Let $\varpi\in K$ satisfy $0<|\varpi|<1$, and consider the compact
annulus and its two closed subannuli
\[
 \begin{split}
 B&=\{|\varpi|\leq |T|\leq |\varpi|^{-1}\},\\
 B_-&=\{|\varpi|\leq |T|\leq1\},\qquad
 B_+=\{1\leq |T|\leq |\varpi|^{-1}\}.
 \end{split}                                                    \tag{HU6}
\]
The map $B_-\amalg B_+\to B$ is surjective and quasi-\'etale.
The unit Gauss point is in the ordinary interior of $B$ and on the
analytic boundary of both displayed source domains.  Thus the
shortcut asserting an ordinary interior lift for every point of a
surjective quasi-\'etale compact cover is \texttt{FALSE}.  This is
the already frozen R19 UF5 example, not a new counterexample family.
It does not contradict HU2: $B_-,B_+$ themselves are a finite
target G-cover, and scalar PL presentations glue on that cover.
The associated R19 UF6--UF8 current example explains why discarding
the boundary is nevertheless invalid for arbitrary-current action.

The following conclusions are deliberately \emph{not} part of HU5.
\begin{enumerate}
\item A finite continuous scalar with ordinary $D[h]=0$ has not,
by the definition alone, been identified with a classically
pluriharmonic scalar.  Any use of HU5 for KC harmonic witnesses must
pin a valid comparison in exactly that scalar class.  No general
positivity-to-regularization implication is inferred from the
abstracts of the cited papers.
\item Geometric real PL on a G-cover is not yet a single global
unit-log expression on a prescribed ordinary domain.  A
residue/Picard argument, when applicable, must retain its own exact
hypotheses and ordinary-germ quantifiers.
\item HU2 and HU5 construct neither pullback nor descent for
arbitrary allowed positive currents along the quasi-\'etale map.
They do not localize a previously chosen full carrier polar.
\item No conclusion about mixed KC harmonic generators with
arbitrary companion potentials, the full carrier-CLN seminorm, or
the actual signed PL relation ideal follows just from scalar
recognition.  No ordinary-current equality is promoted to an
arbitrary action on GK delta tests.
\end{enumerate}

Accordingly the original same-germ KC12 implication and the
variable-PS/CTO gates retain status \texttt{OPEN}.  The precise
topological obstruction to the naive three-result concatenation is
repaired for geometric scalar PL descent by HU2; it is not repaired
for arbitrary-current descent by the same argument.
